% Minimal retained preamble for lane-p2 cards taken from
% sources/logic-fl02-conservation-of-ramsey-s-theorem-f/pinned-original-source.tex
% (Conservation of Ramsey's theorem for pairs, arXiv:2402.11616v3).
%
% Copied verbatim from the author's own file:
%   line 2      : \documentclass{article}
%   lines 16-25 : the author's theorem declarations
%   lines 46-88 : the author's macro definitions for the reverse-mathematics principles
%   line 91     : \card
%   lines 104-107: the natural sums and products, which rely on mathabx (source line 103)
%   line 116    : \qt
%
% Omitted author preamble lines and why (nothing was renamed or redefined):
%   line 3 (graphicx), 7 (url), 9 (cancel): no figure, url or cancel use in the retained spans;
%   line 6 (cleveref): the assembler adds its own cleveref machinery, a second load clashes;
%   line 8 (\usepackage[utf8x]{inputenc}): inputenc with the utf8x option fights the forced
%       xelatex engine;
%   lines 28-37 (the rep@theorem machinery and \newreptheorem{maintheorem}): the retained spans
%       use no rep-theorem environment;
%   lines 40-43 (xcolor and the three author comment macros): not used by the retained spans;
%   lines 93-99 (the LS2 font encoding block for \smashtimes based on stix): the host TeX Live has
%       no stix directory and the author himself never uses the result -- line 105 maps \ntimes to
%       \dottimes, and line 110 shows \smashtimes already commented out in the source.
% mathabx is retained because \dotplus, \dottimes and \dotdiv come from it.
\documentclass{article}
\usepackage{amsmath,amssymb,amsthm}
\usepackage{mathabx}
\usepackage{cancel}
\usepackage{cleveref}

\newtheorem{theorem}{Theorem}
\numberwithin{theorem}{section}
\newtheorem{maintheorem}[theorem]{Main Theorem}
\newtheorem{lemma}[theorem]{Lemma}
\newtheorem{question}[theorem]{Question}
\newtheorem{proposition}[theorem]{Proposition}
\newtheorem{statement}[theorem]{Statement}
\newtheorem{remark}[theorem]{Remark}
\newtheorem{definition}[theorem]{Definition}
\newtheorem{corollary}[theorem]{Corollary}

\newcommand{\RCA}{\mathsf{RCA}}
\newcommand{\WKL}{\mathsf{WKL}}
\newcommand{\SRT}{\mathsf{SRT}}
\newcommand{\EM}{\mathsf{EM}}
\newcommand{\SEM}{\mathsf{SEM}}
\newcommand{\ADS}{\mathsf{ADS}}
\newcommand{\SCAC}{\mathsf{SCAC}}
\newcommand{\SADS}{\mathsf{SADS}}
\newcommand{\COH}{\mathsf{COH}}
\newcommand{\ACA}{\mathsf{ACA}}
\newcommand{\qvdash}{\operatorname{{?}{\vdash}}}
\newcommand{\nqvdash}{\operatorname{{?}{\nvdash}}}
\newcommand{\BSig}{\mathsf{B}\Sigma^0}
\newcommand{\ISig}{\mathsf{I}\Sigma^0}
\newcommand{\LSig}{\mathsf{L}\Sigma^0}
\newcommand{\BPi}{\mathsf{B}\Pi^0}
\newcommand{\IPi}{\mathsf{I}\Pi^0}
\newcommand{\LPi}{\mathsf{L}\Pi^0}
\newcommand{\IDel}{\mathsf{I}\Delta^0}
\newcommand{\BCDe}{\mathsf{BC}\Delta^0}
\newcommand{\WF}{\mathsf{WF}}
\newcommand{\RT}{\mathsf{RT}}
\newcommand{\BME}{\mathsf{BME}}
\newcommand{\TT}{\mathsf{TT}}
\renewcommand{\O}{\mathcal{O}}
\renewcommand{\L}{\mathcal{L}}
\newcommand{\C}{\mathcal{C}}
\newcommand{\R}{\mathcal{R}}
\renewcommand{\S}{\mathcal{S}}
\newcommand{\T}{\mathcal{T}}
\renewcommand{\P}{\mathcal{P}}
\newcommand{\M}{\mathcal{M}}
\newcommand{\Nc}{\mathcal{N}}
\renewcommand{\P}{\mathcal{P}}
\newcommand{\Q}{\mathcal{Q}}
\newcommand{\NN}{\mathbb{N}}
\newcommand{\PP}{\mathbb{P}}
\newcommand{\uh}[0]{{\upharpoonright}}
\newcommand{\card}{\operatorname{card}}
\newcommand{\nplus}{\dotplus}
\newcommand{\ntimes}{\dottimes}
\newcommand{\nminus}{\dotdiv}
\newcommand{\nsum}{\dot{\sum}}
\def\qt#1{``#1''}
\title{A tree condition forces the Shore-blocking functional to be partial}
\author{Quentin Le Houérou, Ludovic Levy Patey, Keita Yokoyama}
\date{arxiv:2402.11616v3; CC BY 4.0}
\begin{document}
\maketitle
\noindent\textit{Editorial note.} The counted unit is the article's own Lemma lem:wkl-wf-forcing-wf at source lines 683-685 together with the proof environment that begins two lines after the statement ends (687-709); the pairing is the ordinary adjacent one.\par\medskip
\noindent\textit{Editorial note.} The retained context at source lines 411-413 and 426-431 is the statement of the two earlier lemmas that the counted proof invokes through the author Cref calls, namely the transfer of well-foundedness to the natural product and the Shore-blocking functional; both are retained as prerequisites, their proofs are not retained. The retained context at source line 663 is the author explanation of the forcing relation for a decided initial segment of the jump.\par\medskip
\noindent\textit{Editorial note.} The retained preamble reproduces verbatim the author's document class at source line 2, his theorem declarations at source lines 16-25, his macro definitions at source lines 46-88 and 91, his natural-operation macros at source lines 104-107 and his quotation macro at source line 116, and it retains mathabx from source line 103 and the cancel package from source line 9, the latter because the retained proof prints a negated satisfaction sign. The LS2 font-encoding block at source lines 93-99 is not retained because the host TeX Live has no stix directory and the author never uses that symbol. The inputenc load with the utf8x option at source line 8 is not retained because it fights the forced xelatex engine, and cleveref is loaded by the retained preamble itself because the assembler only adds that package when it generates navigation stubs.\par\medskip
\noindent\textit{Editorial note.} No retained line contains a citation, so the material carries no bibliography; the cross-references of the retained proof are all resolved inside the retained context.\par\medskip
\noindent\textit{Editorial note.} The article's conservation theorems, its COH lemmas and its SEM section lemmas are not counted here and are not used as premises.\par\medskip
\noindent\textit{Allowed standard prerequisites (editorial).}\begin{itemize}
\item Tree conditions of the forcing used for WKL. In this section a condition is a pair (sigma, T) of a finite binary string and an M-infinite tree, and an extension (tau, T\_2) of (sigma, T\_1) requires tau to extend sigma and T\_2 to be contained in T\_1, so that the conditions form a forcing of tree type; the requirements of this forcing are listed by the author at source lines 640-643.. Fixes the objects of the counted lemma, namely the condition (sigma, T\_1), the extension (sigma, T\_2) it produces and the generics G and Y.
\item The Shore-blocking functional Gamma\_k and well-foundedness of the natural product. Gamma\_k is the Turing functional constructed by the author in the Shore-blocking lemma, which turns a descending sequence below an ordinal into one below the natural product with k, and the article proves separately that well-foundedness of an ordinal transfers to its natural product with any coefficient; both statements are retained with this card at source lines 426-431 and 411-413.. The counted lemma is stated for exactly this functional and its proof derives a contradiction from the hypothesis that Gamma\_k is total, using the transfer of well-foundedness to the natural product.
\item Definability classes and M-finiteness in the model. Sets defined by primitive recursive predicates in a real Y are available in the model, and the article uses throughout that M-infinite trees can be cut down by such predicates while staying M-infinite when the corresponding set is M-infinite.. In the first case the proof cuts the tree T\_1 down by the primitive Y-recursive predicate that the functional diverges at a fixed n.
\end{itemize}\medskip
\begin{lemma}\label[lemma]{lem:wf-product}
Let~$\M = (M,S) \models \RCA_0 + \WF(\alpha)$ for $\alpha$ an ordinal less than $\epsilon_0$. Then $\M \models \WF(\alpha \ntimes k)$ for every $k \in M$.
\end{lemma}

\begin{lemma}\label[lemma]{lem:shore-blocking-wf}
Consider a model $\M = (M,S)$ and $\alpha \in M$ an ordinal less than $\epsilon_0$.
For every~$k \in M$, there is a Turing functional $\Gamma_k$ such that, letting~$\beta \in \alpha$ be the largest ordinal less than $\alpha$ of index at most~$k$, for every~$X \in 2^M$ such that $\M[X] \models \RCA_0$, if there is some~$e < k$ such that $\Phi^X_e$ is an $M$-infinite decreasing sequence of elements smaller than~$\beta$, then $\Gamma^X_k$ is an $M$-infinite decreasing sequence of elements smaller than~$\beta \ntimes k$.

Moreover, an index of~$\Gamma_k$ can be found computably in~$k$.
\end{lemma}

In what follows, $(\tau, U) \Vdash (G \oplus Y)' \uh_k = \sigma'$ means that for every~$e < |\sigma'|$, if $\sigma'(e) = 1$ then $(\tau, U) \Vdash \Phi^{G \oplus Y}_e(e)\downarrow$ and if $\sigma'(e) = 0$, then $(\tau, U) \Vdash \Phi^{G \oplus Y}_e(e)\uparrow$. Thus, this relation is $\Pi^0_1(Y)$.

\begin{lemma}\label[lemma]{lem:wkl-wf-forcing-wf}
Let $(\sigma, T_1)$ be a condition. For every $k \in M$, letting~$\Gamma_k$ be the functional of \Cref{lem:shore-blocking-wf}, there is an extension $(\sigma, T_2) \leq (\sigma,T_1)$ forcing $\Gamma^{G \oplus Y}_k$ to be partial.
\end{lemma}

\begin{proof}
By twisting $\Gamma_k$, we can assume that for every~$e, n \in M$, if $\Gamma^{\sigma \oplus Y}_k(n)\downarrow$, then
\begin{itemize}
    \item $\Gamma^{\sigma \oplus Y}_k(m)\downarrow$ for every~$m < n$ ;
    \item $\Gamma^{\sigma \oplus Y}_k(0), \Gamma^{\sigma \oplus Y}_k(1), \dots, \Gamma^{\sigma \oplus Y}_k(n)$ is a strictly $<_{\epsilon_0}$-decreasing sequence of ordinals smaller than~$\beta \ntimes k$.
\end{itemize}
We have two cases.

Case 1: there exists some~$n \in M$ such that for every~$\ell \in M$, there is some~$\tau \in T_1$ of length~$\ell$ such that $\Gamma^{\tau \oplus Y}_k(n)\uparrow$. Let~$T_2 = \{ \tau \in T_1 : \Gamma^{\tau \oplus Y}_k(n)\uparrow \}$. Note that the set $T_2$ is primitive $Y$-recursive, as the set~$T_1$ and the predicate $\Gamma^{\sigma \oplus Y}_k(n)\uparrow$ are primitive $Y$-recursive. The set~$T_2$ is closed-downward, and by assumption, is $M$-infinite. Then $(\sigma, T_2) \leq (\sigma, T_1)$. Moreover, by construction, $(\sigma, T_2) \Vdash \Gamma_k^{G \oplus Y}(n)\uparrow$.


Case 2: for every~$n \in M$, there is some~$\ell_n \in M$ such that 
for every~$\tau \in T$ of length~$\ell_n$, $\Gamma^{\tau \oplus Y}_k(n)\downarrow$.

Then, for every~$n$, let $\alpha_n = \max\ \{ \Gamma^{\tau \oplus Y}_k(n) : \tau \in T_1 \wedge |\tau| = \ell_n \}$.
We claim that for every~$n \in M$, $\alpha_{n+1} <_{\epsilon_0} \alpha_n$.
Indeed, for every~$\tau \in T_1$ such that $|\tau| = \ell_{n+1}$, $\Gamma^{\tau \oplus Y}_k(n+1) <_{\epsilon_0} \Gamma^{(\tau \uh \ell_n) \oplus Y}_k$, so
$$
\max\ \{ \Gamma^{\tau \oplus Y}_k(n+1) : \tau \in T_1 \wedge |\tau| = \ell_{n+1} \}
<_{\epsilon_0} \max\ \{ \Gamma^{\tau \oplus Y}_k(n) : \tau \in T_1 \wedge |\tau| = \ell_n \}
$$
So $\M \mathbin{\cancel{\models}} \WF(\alpha \ntimes k)$. However, by \Cref{lem:wf-product}, since $\M \models \BSig_2 + \WF(\alpha)$, $\M \models \WF(\alpha \ntimes k)$. Contradiction.
\end{proof}

\end{document}
